\section{Técnicas Básicas}

Esta sección detalla las herramientas fundamentales para controlar formato, tipografía y estructura en \LaTeX.

\subsection{Párrafos y Separación de Texto}

Los párrafos se crean automáticamente al dejar una línea en blanco entre bloques de texto \cite{latex_project}. Por ejemplo:

\begin{verbatim}
Este es el primer párrafo.

Este es el segundo párrafo.
\end{verbatim}

\begin{tcolorbox}[title=Salida,colback=white!10,colframe=black,fonttitle=\bfseries]
Este es el primer párrafo.

Este es el segundo párrafo.
\end{tcolorbox}

También se pueden usar \texttt{\textbackslash\textbackslash} para saltos de línea forzosos dentro de un mismo párrafo:

\begin{verbatim}
Esta es la primera línea.\\
Continuamos en la siguiente línea.
\end{verbatim}

\begin{tcolorbox}[title=Salida,colback=white!10,colframe=black,fonttitle=\bfseries]
Esta es la primera línea.\\
Continuamos en la segunda línea.
\end{tcolorbox}

Tanto la sangría como el espacio entre párrafos se pueden configurar con \texttt{\textbackslash setlength}:

\begin{verbatim}
\setlength{\parindent}{1.25cm} 
\setlength{\parskip}{1em}
\end{verbatim}

El primer comando establece la sangría de párrafo a 1.25 cm, mientras que el segundo agrega un espacio vertical de 1 em entre párrafos.

\subsection{Efectos de Letra y Tipografía}

\LaTeX ofrece comandos para modificar el estilo de texto, como negrita, itálica y monoespaciado \cite{latex_fontguide}. A continuación se muestran los comandos básicos:

\begin{itemize}
    \item \textbf{Negrita}: \texttt{\textbackslash textbf\{texto\}}
    \item \textit{Itálica}: \texttt{\textbackslash textit\{texto\}}
    \item \texttt{Monoespaciado}: \texttt{\textbackslash texttt\{texto\}}
    \item \textsf{Sans-serif}: \texttt{\textbackslash textsf\{texto\}}
    \item \textsc{Pequeñas mayúsculas}: \texttt{\textbackslash textsc\{texto\}}
    \item \emph{Énfasis contextual}: \texttt{\textbackslash emph\{texto\}}
\end{itemize}

Asimismo, se pueden cambiar las familias de fuentes completas:

\begin{itemize}
    \item \textbf{Roman}: \texttt{\textbackslash rmfamily\{texto\}}
    \item \textbf{Sans-serif}: \texttt{\textbackslash sffamily\{texto\}}
    \item \textbf{Monoespaciado}: \texttt{\textbackslash ttfamily\{texto\}}
\end{itemize}

\subsection{Tildes y Caracteres Especiales}

Para escribir caracteres acentuados en español, se recomienda usar UTF-8 con \texttt{\textbackslash usepackage[utf8]\{inputenc\}}. Esto permite escribir directamente:

\begin{verbatim}
Café, mañana, pingüino
\end{verbatim}

Sin embargo, también es posible usar comandos de acentuación y obtener el mismo resultado:

\begin{verbatim}
Caf\'{e}, ma\~{n}ana, ping\"{u}ino
\end{verbatim}

\subsection{Títulos, Subtítulos y Estructura}

\LaTeX proporciona comandos para organizar el documento en secciones, subsecciones y niveles inferiores \cite{latex_project}. La sintaxis general es:

\begin{verbatim}
\section{Título principal}           % Nivel 1
\subsection{Subtítulo}               % Nivel 2
\subsubsection{Sub-subtítulo}        % Nivel 3
\paragraph{Título párrafo}           % Nivel 4
\subparagraph{Sub-párrafo}           % Nivel 5
\end{verbatim}

\subsection{Encabezados, Pies de Página y Marcas de Agua}

Se pueden configurar encabezados y pies de página con fancyhdr \cite{fancyhdr}:

\begin{verbatim}
\usepackage{fancyhdr}
\pagestyle{fancy}
\fancyhf{}
\fancyhead[L]{Trabajo \LaTeX{} 2026} 
\fancyhead[R]{\thepage} 
\fancyfoot[C]{\small Costa Rica}
\renewcommand{\headrulewidth}{0.4pt}
\end{verbatim}

También se puede usar draftwatermark \cite{draftwatermark} para agregar marcas de agua:

\begin{verbatim}
\usepackage{draftwatermark}
\SetWatermarkText{\textbf{DRAFT}}
\SetWatermarkScale{0.3}
\SetWatermarkColor{rgb}{0.8,0.8,0.8}
\end{verbatim}

\subsection{Referencias y Citas}

Para citar fuentes, se recomienda usar BibLaTeX con un archivo .bib \cite{biblatex}. Ejemplo de cita:

\begin{verbatim}
\cite{latex}
\end{verbatim}

La referencia se define en un archivo .bib de la siguiente forma:

\begin{verbatim}
@online{latex,
  author = {LaTeX Project Team},
  title = {LaTeX2e User's Guide},
  year = {2025},
  url = {https://www.latex-project.org},
  urldate = {2026-04-03}
}
\end{verbatim}

\subsection{Saltos de Página y Control de Flujo}

Los siguientes comandos controlan el flujo del documento:

\begin{itemize}
    \item \textbf{\textbackslash newpage}: Salto de página forzoso
    \item \textbf{\textbackslash clearpage}: Salto de página con colocación de figuras pendientes (tablas, gráficos)
    \item \textbf{\textbackslash pagebreak}: Sugerencia suave para salto de página
    \item \textbf{\textbackslash newcolumn}: Salto de columna dentro de multicols
    \item \textbf{\textbackslash columnbreak}: Salto de columna forzoso
\end{itemize}

\subsection{Columnas de Página con multicol}

Se puede dividir el texto en múltiples columnas usando el paquete multicol \cite{multicol}. Esto es útil para formatos de revista o boletines:

\begin{verbatim}
\usepackage{multicol}
\begin{multicols}{2}
Columna izquierda con texto largo...

Columna derecha continúa automáticamente.
\end{multicols}
\end{verbatim}

\begin{tcolorbox}[title=Salida,colback=white!10,colframe=black,fonttitle=\bfseries]
\begin{multicols}{2}
Columna izquierda con texto largo...

Columna derecha continúa automáticamente.
\end{multicols}
\end{tcolorbox}